\documentclass[11pt]{article}

\usepackage[margin=1in]{geometry}
\usepackage{amsmath,amssymb,amsthm}
\usepackage{booktabs}
\usepackage[hidelinks]{hyperref}

\newtheorem{theorem}{Theorem}[section]
\newtheorem{proposition}[theorem]{Proposition}
\newtheorem{lemma}[theorem]{Lemma}
\newtheorem{corollary}[theorem]{Corollary}
\theoremstyle{definition}
\newtheorem{definition}[theorem]{Definition}
\newtheorem{remark}[theorem]{Remark}

\newcommand{\FD}{\mathrm{FD}}
\newcommand{\FDz}{\mathrm{FD}_0}
\newcommand{\Maj}{\mathrm{Maj}}

\title{\textbf{Disproof of a Logarithmic Central-Correction Conjecture\\
for the Sperner Number of the Free Distributive Lattice}}
\author{Formalized Resolution of Conjecture 00000008532}
\date{}

\begin{document}
\maketitle

\begin{abstract}
The \emph{Sperner number} (or \emph{width}) $w(L)$ of a poset $L$ is the maximum cardinality of an antichain in $L$. Conjecture~00000008532 asserts that for the free distributive lattice on $n$ generators, $w(\FD(n))$ equals $\binom{n}{\lfloor n/2 \rfloor} + \Delta(n)$ where the central correction $\Delta(n)$ is of logarithmic order and has strictly positive sign for all $n$. We disprove both assertions under both standard conventions for the free distributive lattice: the bounded free distributive lattice $\FD(n)$ (isomorphic to the lattice of all monotone Boolean functions $\{0,1\}^n \to \{0,1\}$, of cardinality the $n$-th Dedekind number $M(n)$) and the unbounded free distributive lattice $\FDz(n)$ (without empty join or empty meet, of cardinality $M(n) - 2$ for $n \ge 1$). First, we prove that $w(\FD(n)) = w(\FDz(n))$ for every $n \ge 1$ and compute the exact values $w(\FD(1)) = 1 = \binom{1}{0}$ and $w(\FD(2)) = 2 = \binom{2}{1}$, so $\Delta(1) = \Delta(2) = 0$ (and $\Delta(0) = 0$ in $\FD(0)$, $\Delta_0(0) = -1$ in $\FDz(0)$), refuting strict positivity. Second, by embedding the middle layer of subsets of the middle Boolean slice $\binom{[n]}{\lfloor n/2 \rfloor}$ as monotone Boolean functions, we construct an explicit antichain in $\FD(n)$ and $\FDz(n)$ of cardinality
\[
  \binom{N_n}{\lfloor N_n / 2 \rfloor} \;\ge\; \frac{2^{N_n}}{N_n + 1} \;\ge\; 2^{2^{\lfloor n/2 \rfloor - 1}}
  \qquad (n \ge 2),
  \qquad N_n = \binom{n}{\lfloor n/2 \rfloor}.
\]
Consequently, $\Delta(n) > 0$ holds if and only if $n \ge 3$, and $\Delta(n)$ grows doubly exponentially in $n$---exceeding $C B^n + D$ for every base $B > 0$ and constants $C, D$, and in particular refuting $\Delta(n) = O(\log n)$. Combined with Kleitman's upper bound on $M(n)$, it is $\log_2 w(\FD(n))$, rather than $w(\FD(n))$ itself, that has asymptotic order $\binom{n}{\lfloor n/2 \rfloor}$. All core results are machine-checked from first principles in Core Lean~4.19.0 without external libraries.
\end{abstract}

\section{Introduction}

Recall that an \emph{antichain} in a partially ordered set $(P, \le)$ is a subset $A \subseteq P$ of pairwise incomparable elements, i.e., $x \le y \implies x = y$ for all $x, y \in A$. The \emph{Sperner number} (or \emph{width}) of $P$, denoted $w(P)$, is the maximum cardinality of an antichain in $P$. By Sperner's theorem~\cite{Sperner1928}, the Boolean lattice $B_n = (2^{[n]}, \subseteq)$ of all subsets of $[n] = \{1, \dots, n\}$ has Sperner number
\[
  w(B_n) = \binom{n}{\lfloor n/2 \rfloor},
\]
attained by the middle layer $\binom{[n]}{\lfloor n/2 \rfloor}$.

Conjecture~00000008532 proposes an analogous formula for the free distributive lattice on $n$ generators:
\begin{quote}
\emph{Definition: The Sperner number of a lattice is the size of its largest antichain. Conjecture: The Sperner number of the free distributive lattice is an explicit central correction of $\binom{n}{\lfloor n/2 \rfloor}$; the correction is of logarithmic order with always positive sign.}
\end{quote}

In the literature on lattice theory and Dedekind's problem~\cite{Dedekind1897,Birkhoff1937,Kleitman1969,Engel1997}, the phrase ``free distributive lattice on $n$ generators'' is used in two closely related senses:
\begin{enumerate}
  \item The \emph{bounded} free $(0,1)$-distributive lattice $\FD(n)$, allowing empty join $\bot$ and empty meet $\top$. Concretely, $\FD(n)$ is isomorphic to the pointwise-ordered lattice of all monotone Boolean functions $f : \{0,1\}^n \to \{0,1\}$, whose cardinality $|\FD(n)| = M(n)$ is the $n$-th Dedekind number ($2, 3, 6, 20, 168, 7581, \dots$ for $n = 0, 1, 2, 3, 4, 5$).
  \item The \emph{unbounded} free distributive lattice $\FDz(n)$, generated by $x_1, \dots, x_n$ using only non-empty finite joins and meets. Under the monotone Boolean representation, $\FDz(n)$ consists of all monotone functions $f : \{0,1\}^n \to \{0,1\}$ satisfying $f(0,\dots,0) = 0$ and $f(1,\dots,1) = 1$, with $|\FDz(0)| = 0$ and $|\FDz(n)| = M(n) - 2$ for $n \ge 1$ ($1, 4, 18, 166, 7579, \dots$).
\end{enumerate}

At first glance, comparing $w(\FD(n))$ with $\binom{n}{\lfloor n/2 \rfloor}$ has a natural motivation: for each subset $S \subseteq [n]$ of size $\lfloor n/2 \rfloor$, the elementary conjunction $\bigwedge_{i \in S} x_i$ defines an element of $\FD(n)$, and these $\binom{n}{\lfloor n/2 \rfloor}$ conjunctions are pairwise incomparable. Thus $w(\FD(n)) \ge \binom{n}{\lfloor n/2 \rfloor}$ holds for every $n \ge 0$, so the signed central correction
\[
  \Delta(n) := w(\FD(n)) - \binom{n}{\lfloor n/2 \rfloor}, \qquad
  \Delta_0(n) := w(\FDz(n)) - \binom{n}{\lfloor n/2 \rfloor}
\]
satisfies $\Delta(n) \ge 0$ for all $n \ge 0$ (and $\Delta_0(n) = \Delta(n) \ge 0$ for all $n \ge 1$). Moreover, at $n = 3$ one finds $w(\FD(3)) = 4 = \binom{3}{1} + 1$, so $\Delta(3) = 1$, which superficially resembles a small positive correction.

In this paper we prove that \textbf{Conjecture~00000008532 is false in both claims and under both conventions $\FD(n)$ and $\FDz(n)$}. Our main results are summarized in the following two theorems.

\begin{theorem}[Invariance, Exact Small-$n$ Widths, and Sign Classification]\label{thm:small-and-sign}
For every integer $n \ge 1$, the bounded and unbounded free distributive lattices have identical Sperner numbers:
\[
  w(\FD(n)) = w(\FDz(n)) \qquad \text{and} \qquad \Delta(n) = \Delta_0(n).
\]
For $0 \le n \le 5$, the exact Sperner numbers and central corrections are:
\begin{itemize}
  \item $n = 0$: $w(\FD(0)) = 1 = \binom{0}{0}$ ($\Delta(0) = 0$), while $w(\FDz(0)) = 0$ ($\Delta_0(0) = -1 < 0$);
  \item $n = 1$: $w(\FD(1)) = w(\FDz(1)) = 1 = \binom{1}{0}$, so $\Delta(1) = \Delta_0(1) = 0$;
  \item $n = 2$: $w(\FD(2)) = w(\FDz(2)) = 2 = \binom{2}{1}$, so $\Delta(2) = \Delta_0(2) = 0$;
  \item $n = 3$: $w(\FD(3)) = w(\FDz(3)) = 4$, while $\binom{3}{1} = 3$, so $\Delta(3) = \Delta_0(3) = 1$;
  \item $n = 4$: $w(\FD(4)) = w(\FDz(4)) = 24$, while $\binom{4}{2} = 6$, so $\Delta(4) = \Delta_0(4) = 18$;
  \item $n = 5$: $w(\FD(5)) = w(\FDz(5)) = 621$, while $\binom{5}{2} = 10$, so $\Delta(5) = \Delta_0(5) = 611$.
\end{itemize}
In general, $\Delta(n) \ge 0$ for all $n \ge 0$, and $\Delta(n) > 0 \iff n \ge 3$ (likewise $\Delta_0(n) > 0 \iff n \ge 3$ for all $n \ge 1$). In particular, the correction does \emph{not} have always-positive sign.
\end{theorem}

\begin{theorem}[Doubly-Exponential Lower Bound and Refutation of Logarithmic Order]\label{thm:double-exp}
Let $N_n = \binom{n}{\lfloor n/2 \rfloor}$. For every $n \ge 0$ (and in $\FDz(n)$ for every $n \ge 1$), the Sperner number satisfies
\begin{equation}\label{eq:iter-cb}
  w(\FD(n)) \;\ge\; \binom{N_n}{\lfloor N_n / 2 \rfloor} \;\ge\; \frac{2^{N_n}}{N_n + 1},
\end{equation}
and for every $n \ge 2$,
\begin{equation}\label{eq:clean-dexp}
  w(\FD(n)) = w(\FDz(n)) \;\ge\; 2^{2^{\lfloor n/2 \rfloor - 1}}.
\end{equation}
Consequently, for every threshold $n_0 \in \mathbb{N}$, base $B \in \mathbb{N}$, and constants $C, D \in \mathbb{N}$, there exists $n \ge n_0$ such that
\[
  \Delta(n) = \Delta_0(n) > C B^n + D.
\]
In particular, $\Delta(n)$ is neither $O(\log n)$, nor $O(n^d)$ for any $d$, nor $O(B^n)$ for any $B$.
\end{theorem}

\section{Free Distributive Lattices and Width Invariance}

Let $\{0,1\}^n$ carry the coordinatewise order $x \le y \iff (\forall i \in [n],\, x_i \le y_i)$, with minimum $\mathbf{0} = (0,\dots,0)$, maximum $\mathbf{1} = (1,\dots,1)$, and Hamming weight $|x| = \sum_{i=1}^n x_i$. If $x \le y$, then $|x| \le |y|$, with equality $|x| = |y|$ if and only if $x = y$.

\begin{definition}[Monotone Boolean Models of $\FD(n)$ and $\FDz(n)$]\label{def:fd}
For $n \ge 0$:
\begin{enumerate}
  \item $\FD(n)$ is the set of all monotone functions $f : \{0,1\}^n \to \{0,1\}$ (meaning $x \le y \implies f(x) \le f(y)$), ordered pointwise by $f \le g \iff (\forall x \in \{0,1\}^n,\, f(x) \le g(x))$, with pointwise join $(f \vee g)(x) = \max(f(x), g(x))$, pointwise meet $(f \wedge g)(x) = \min(f(x), g(x))$, bottom $\bot(x) = 0$, top $\top(x) = 1$, and $n$ generators $g_i(x) = x_i$ ($1 \le i \le n$).
  \item $\FDz(n) = \{f \in \FD(n) : f(\mathbf{0}) = 0 \text{ and } f(\mathbf{1}) = 1\}$, with the induced order, join, and meet.
\end{enumerate}
\end{definition}

Let us briefly recall why Definition~\ref{def:fd} realizes the categorical free objects (we also formalize the universal property of $\FD(n)$ directly in Lean~4).

\begin{proposition}[Universal Property and DNF Representation]\label{prop:universal}
For every $n \ge 0$ and every bounded distributive lattice $(L, \le, \vee, \wedge, 0_L, 1_L)$, any map $v : [n] \to L$ extends to a unique bounded lattice homomorphism $\Phi_v : \FD(n) \to L$ satisfying $\Phi_v(g_i) = v(i)$ for all $i \in [n]$. Moreover, for $n \ge 1$, $\FDz(n)$ is precisely the sublattice of $\FD(n)$ generated by $\{g_1, \dots, g_n\}$ under binary $\vee$ and $\wedge$.
\end{proposition}

\begin{proof}
We proceed by induction on $n$ via Shannon decomposition. For $n = 0$, $\FD(0) = \{\bot, \top\}$ and $\Phi_v(\bot) = 0_L$, $\Phi_v(\top) = 1_L$ is uniquely forced. For $n + 1$, every $f \in \FD(n + 1)$ restricts to cofactors $f_0, f_1 \in \FD(n)$ given by $f_b(x_2,\dots,x_{n+1}) = f(b, x_2, \dots, x_{n+1})$ with $f_0 \le f_1$, and satisfies the identity
\[
  f = \tilde{f}_0 \vee (g_1 \wedge \tilde{f}_1)
\]
in $\FD(n + 1)$, where $\tilde{h}(x_1, x_2, \dots, x_{n+1}) = h(x_2, \dots, x_{n+1})$. Defining $\Phi_v(f) := \Phi_{v'}(f_0) \vee (v(1) \wedge \Phi_{v'}(f_1))$ where $v'(j) = v(j + 1)$, monotonicity of $\Phi_{v'}$ gives $a_0 := \Phi_{v'}(f_0) \le \Phi_{v'}(f_1) =: a_1$. Using distributivity in $L$, one checks immediately that $\Phi_v$ preserves $\bot, \top, \vee$ and maps $g_i \mapsto v(i)$; for meets, if $a_0 \le a_1$ and $b_0 \le b_1$ in $L$, expanding $(a_0 \vee (v(1) \wedge a_1)) \wedge (b_0 \vee (v(1) \wedge b_1))$ yields four joinands, of which $a_0 \wedge v(1) \wedge b_1$ and $v(1) \wedge a_1 \wedge b_0$ are both bounded above by $v(1) \wedge a_1 \wedge b_1$ and absorbed into it, proving $\Phi_v(f \wedge h) = \Phi_v(f) \wedge \Phi_v(h)$. Uniqueness follows from the Shannon identity in $\FD(n + 1)$ and the induction hypothesis applied to $h \mapsto \Phi(\tilde{h})$.

Finally, for $n \ge 1$, every generator $g_i$ lies in $\FDz(n)$, and $\FDz(n)$ is closed under binary $\vee$ and $\wedge$. Conversely, if $f \in \FDz(n)$, the set $M_f$ of minimal $a \in \{0,1\}^n$ with $f(a) = 1$ is non-empty (since $f(\mathbf{1}) = 1$) and excludes $\mathbf{0}$ (since $f(\mathbf{0}) = 0$), so $f = \bigvee_{a \in M_f} \bigwedge_{i : a_i = 1} g_i$ is a non-empty join of non-empty meets of generators.
\end{proof}

\begin{proposition}[Width Invariance]\label{prop:invariance}
For every integer $n \ge 1$, we have $\FD(n) = \{\bot\} \sqcup \FDz(n) \sqcup \{\top\}$, and consequently
\[
  w(\FD(n)) = w(\FDz(n)).
\]
\end{proposition}

\begin{proof}
Since $n \ge 1$, $\mathbf{0} \ne \mathbf{1}$ in $\{0,1\}^n$, so $\bot \ne \top$. For any $f \in \FD(n)$, monotonicity between $\mathbf{0} \le \mathbf{1}$ implies $(f(\mathbf{0}), f(\mathbf{1})) \in \{(0,0), (0,1), (1,1)\}$. If $f(\mathbf{1}) = 0$, then $x \le \mathbf{1}$ forces $f(x) = 0$ for all $x$, so $f = \bot$; if $f(\mathbf{0}) = 1$, then $\mathbf{0} \le x$ forces $f(x) = 1$ for all $x$, so $f = \top$; the remaining case $(0,1)$ defines $\FDz(n)$.

Every antichain in $\FDz(n)$ is an antichain in $\FD(n)$, so $w(\FDz(n)) \le w(\FD(n))$. Conversely, $\FDz(n)$ is non-empty for $n \ge 1$ (it contains $g_1$), so $w(\FDz(n)) \ge 1$. Let $A \subseteq \FD(n)$ be an antichain of maximum size. If $|A| \le 1$, then $|A| \le 1 \le w(\FDz(n))$. If $|A| \ge 2$, then $A$ cannot contain $\bot$ (since $\bot \le h$ for any other $h \in A$) and cannot contain $\top$ (since $h \le \top$ for any other $h \in A$). Hence $A \subseteq \FDz(n)$, giving $w(\FD(n)) = |A| \le w(\FDz(n))$.
\end{proof}

\begin{proposition}[Classic Sperner Conjunction Lower Bound]\label{prop:conj-lb}
For every $n \ge 0$ and $0 \le k \le n$, we have $w(\FD(n)) \ge \binom{n}{k}$. In particular, $w(\FD(n)) \ge \binom{n}{\lfloor n/2 \rfloor}$ and $\Delta(n) \ge 0$ for all $n \ge 0$ (and $\Delta_0(n) = \Delta(n) \ge 0$ for all $n \ge 1$).
\end{proposition}

\begin{proof}
For each $u \in \{0,1\}^n$ with $|u| = k$, define the principal upset $m_u \in \FD(n)$ by $m_u(x) = 1 \iff u \le x$. Notice that $m_u \le m_v$ in $\FD(n)$ if and only if $v \le u$ in $\{0,1\}^n$ (evaluating $m_u(u) = 1$ yields $m_v(u) = 1$, i.e., $v \le u$, and the converse holds by transitivity). If $|u| = |v| = k$ and $m_u \le m_v$, then $v \le u$ and $|v| = |u|$ force $u = v$. Thus $\{m_u : |u| = k\}$ is an antichain in $\FD(n)$ of size $\binom{n}{k}$.
\end{proof}

\section{Exact Small-\texorpdfstring{$n$}{n} Widths and Failure of Always-Positive Sign}

Recall the elementary half of Dilworth's theorem~\cite{Dilworth1950}: if a poset $P$ can be covered by $c$ chains $C_1, \dots, C_c$, then every antichain $A \subseteq P$ can contain at most one element from each $C_j$, so $w(P) \le c$. Exhibiting both an antichain of size $w$ and a partition of $P$ into $w$ chains therefore certifies $w(P) = w$ without appeal to the full min-max theorem.

\begin{proof}[Proof of the Small-$n$ Values in Theorem~\ref{thm:small-and-sign}]
By Proposition~\ref{prop:invariance}, it suffices to analyze $\FD(n)$ for $n \ge 1$ together with the trivial case $n = 0$.
\begin{itemize}
  \item \textbf{$n = 0$:} $\FD(0) = \{\bot < \top\}$ is a single chain of size $2$, so $w(\FD(0)) = 1 = \binom{0}{0}$ and $\Delta(0) = 0$. Meanwhile $\FDz(0) = \emptyset$, so $w(\FDz(0)) = 0$ and $\Delta_0(0) = 0 - 1 = -1 < 0$.
  \item \textbf{$n = 1$:} $\FD(1) = \{\bot < g_1 < \top\}$ is a $3$-element chain (and $\FDz(1) = \{g_1\}$ is a $1$-element chain). Hence $w(\FD(1)) = w(\FDz(1)) = 1 = \binom{1}{0}$, giving $\Delta(1) = \Delta_0(1) = 0$.
  \item \textbf{$n = 2$:} $\FD(2) = \{\bot, g_1 \wedge g_2, g_1, g_2, g_1 \vee g_2, \top\}$ has $6$ elements. Here $\{g_1, g_2\}$ is an antichain of size $2$, while $\FD(2)$ is partitioned into the $2$ chains
  \[
    C_1 = \{\bot < g_1 \wedge g_2 < g_1 < g_1 \vee g_2 < \top\}, \qquad C_2 = \{g_2\}.
  \]
  Thus $w(\FD(2)) = w(\FDz(2)) = 2 = \binom{2}{1}$, giving $\Delta(2) = \Delta_0(2) = 0$.
  \item \textbf{$n = 3$:} $|\FD(3)| = 20$. Let $\Maj_3 = (g_1 \wedge g_2) \vee (g_1 \wedge g_3) \vee (g_2 \wedge g_3)$ be the majority function. Each of the four elements in $A_3 = \{g_1, g_2, g_3, \Maj_3\}$ takes the value $1$ on exactly $4$ vertices of $\{0,1\}^3$, so no two distinct elements of $A_3$ can be comparable (strict inequality $f < h$ forces $|f^{-1}(1)| < |h^{-1}(1)|$). Encoding $f \in \FD(3)$ by the $8$-bit integer $\sum_{x=0}^7 f(x) 2^x$ where bit $i$ of $x$ is $x_{i+1}$, $A_3$ corresponds to $\{170, 204, 232, 240\}$, and $\FD(3)$ partitions into the $4$ explicit chains
  \begin{align*}
    C_1 &= [0, 128, 136, 168, 170, 234, 238, 255], &
    C_2 &= [160, 224, 232, 248, 254], \\
    C_3 &= [192, 200, 204, 236, 252], &
    C_4 &= [240, 250].
  \end{align*}
  Hence $w(\FD(3)) = w(\FDz(3)) = 4$ and $\Delta(3) = 4 - \binom{3}{1} = 1$.
  \item \textbf{$n = 4$ and $n = 5$:} Grading $\FD(n)$ by rank $r(f) = |f^{-1}(1)| \in \{0, 1, \dots, 2^n\}$, every rank level $L_r = \{f \in \FD(n) : |f^{-1}(1)| = r\}$ is an antichain. At $n = 4$ ($|\FD(4)| = 168$), the middle rank $L_8$ has size $24$ (verified in \texttt{SmallN.lean}), and Hopcroft--Karp matching yields an explicit partition of $\FD(4)$ into $24$ chains (certified in \texttt{verify.py}), proving $w(\FD(4)) = w(\FDz(4)) = 24$ and $\Delta(4) = 24 - 6 = 18$. Similarly, at $n = 5$ ($|\FD(5)| = 7581$), $|L_{16}| = 621$ and an explicit $621$-chain partition certifies $w(\FD(5)) = w(\FDz(5)) = 621$ and $\Delta(5) = 621 - 10 = 611$.
\end{itemize}
The cases $n = 1$ and $n = 2$ (as well as $n = 0$) immediately refute the claim that $\Delta(n)$ has always-positive sign; the remaining equivalence $\Delta(n) > 0 \iff n \ge 3$ follows once we establish Theorem~\ref{thm:double-exp} for $n \ge 4$ below.
\end{proof}

\begin{table}[ht]
\centering
\caption{Exact Dedekind numbers, Sperner numbers, signed central corrections, and middle-slice lower bounds for $0 \le n \le 5$.}
\label{tab:small-n}
\smallskip
\begin{tabular}{@{}rrrrrrrr@{}}
\toprule
$n$ & $|\FD(n)|$ & $|\FDz(n)|$ & $N_n = \binom{n}{\lfloor n/2 \rfloor}$ & $w(\FD(n))$ & $w(\FDz(n))$ & $\Delta(n)$ & $\binom{N_n}{\lfloor N_n/2 \rfloor}$ \\
\midrule
$0$ & $2$    & $0$    & $1$  & $1$   & $0$   & $0$   & $1$   \\
$1$ & $3$    & $1$    & $1$  & $1$   & $1$   & $0$   & $1$   \\
$2$ & $6$    & $4$    & $2$  & $2$   & $2$   & $0$   & $2$   \\
$3$ & $20$   & $18$   & $3$  & $4$   & $4$   & $1$   & $3$   \\
$4$ & $168$  & $166$  & $6$  & $24$  & $24$  & $18$  & $20$  \\
$5$ & $7581$ & $7579$ & $10$ & $621$ & $621$ & $611$ & $252$ \\
\bottomrule
\end{tabular}
\end{table}

\section{Doubly-Exponential Growth and Refutation of Logarithmic Order}

We now construct an explicit doubly-exponential antichain in $\FD(n)$ for all $n$. Let $m = \lfloor n/2 \rfloor$ and let
\[
  \mathcal{M}_n = \{x \in \{0,1\}^n : |x| = m\}
\]
denote the middle weight layer of the Boolean cube, of cardinality $|\mathcal{M}_n| = \binom{n}{m} = N_n$.

\begin{lemma}[Middle-Slice Order Embedding]\label{lem:mid-slice}
For every subset $S \subseteq \mathcal{M}_n$, define $f_S : \{0,1\}^n \to \{0,1\}$ by
\[
  f_S(x) = \begin{cases}
    1 & \text{if } |x| > m, \\
    \mathbf{1}_{x \in S} & \text{if } |x| = m, \\
    0 & \text{if } |x| < m.
  \end{cases}
\]
Then $f_S \in \FD(n)$ for every $S \subseteq \mathcal{M}_n$ (and $f_S \in \FDz(n)$ whenever $n \ge 2$). Moreover, for any $S, T \subseteq \mathcal{M}_n$,
\[
  f_S \le f_T \quad \Longleftrightarrow \quad S \subseteq T.
\]
\end{lemma}

\begin{proof}
Suppose $x \le y$ in $\{0,1\}^n$ and $f_S(x) = 1$. Then $|x| \ge m$ by definition of $f_S$. Since $x \le y$, we have $|x| \le |y|$. If $|y| > m$, then $f_S(y) = 1$. Otherwise $|y| = m$, which forces $|x| = |y| = m$ and hence $x = y$, so $f_S(y) = f_S(x) = 1$. Thus $f_S$ is monotone. When $n \ge 2$, we have $1 \le m = \lfloor n/2 \rfloor \le n - 1$, so $|\mathbf{0}| = 0 < m$ gives $f_S(\mathbf{0}) = 0$ and $|\mathbf{1}| = n > m$ gives $f_S(\mathbf{1}) = 1$, placing $f_S$ in $\FDz(n)$.

If $S \subseteq T$, then $\mathbf{1}_{x \in S} \le \mathbf{1}_{x \in T}$ on $\mathcal{M}_n$ while $f_S$ and $f_T$ agree outside $\mathcal{M}_n$, so $f_S \le f_T$. Conversely, if $f_S \le f_T$ and $x \in S \subseteq \mathcal{M}_n$, then $1 = f_S(x) \le f_T(x)$ implies $f_T(x) = 1$, whence $x \in T$.
\end{proof}

\begin{lemma}[Elementary Central Binomial Estimates]\label{lem:binom-est}
For every integer $M \ge 0$:
\begin{enumerate}
  \item $(M + 1)\binom{M}{\lfloor M/2 \rfloor} \ge 2^M$;
  \item $\binom{M}{\lfloor M/2 \rfloor} \ge 2^{\lfloor M/2 \rfloor}$;
  \item if $M \ge 4$, then $\binom{M}{\lfloor M/2 \rfloor} > M$.
\end{enumerate}
\end{lemma}

\begin{proof}
(1) From $(k + 1)\binom{M}{k + 1} = (M - k)\binom{M}{k}$, we have $\binom{M}{k} \le \binom{M}{k + 1}$ for $k + 1 \le \lfloor M/2 \rfloor$, and by symmetry $\binom{M}{k} = \binom{M}{M - k}$ it follows that $\binom{M}{k} \le \binom{M}{\lfloor M/2 \rfloor}$ for all $0 \le k \le M$. Summing over $0 \le k \le M$ yields $2^M = \sum_{k=0}^M \binom{M}{k} \le (M + 1)\binom{M}{\lfloor M/2 \rfloor}$.

(2) For $M \in \{0, 1, 2, 3\}$, direct inspection gives $1 \ge 1$, $1 \ge 1$, $2 \ge 2$, and $3 \ge 2$. For $M \ge 2$, let $k = \lfloor M/2 \rfloor \ge 1$; then $\lfloor (M + 2)/2 \rfloor = k + 1$, and applying Pascal's identity twice gives
\[
  \binom{M + 2}{k + 1} = \binom{M}{k + 1} + 2\binom{M}{k} + \binom{M}{k - 1} \;\ge\; 2\binom{M}{k},
\]
so $\binom{M}{\lfloor M/2 \rfloor} \ge 2^{\lfloor M/2 \rfloor}$ follows by induction on $M$ with step $2$.

(3) If $M \ge 4$, then $2 \le \lfloor M/2 \rfloor \le M - 2$ and $M - 1 \ge 3$, so $\binom{M}{\lfloor M/2 \rfloor} \ge \binom{M}{2} = \frac{M(M - 1)}{2} \ge \frac{3M}{2} > M$.
\end{proof}

\begin{proof}[Proof of Theorem~\ref{thm:double-exp} and the Sign Equivalence in Theorem~\ref{thm:small-and-sign}]
Define the middle-slice family
\[
  \mathcal{F}_n = \bigl\{f_S : S \subseteq \mathcal{M}_n,\; |S| = \lfloor N_n / 2 \rfloor\bigr\} \subseteq \FD(n).
\]
By Lemma~\ref{lem:mid-slice}, $S \mapsto f_S$ is injective and order-reflecting: if $f_S \le f_T$ for $f_S, f_T \in \mathcal{F}_n$, then $S \subseteq T$ and $|S| = |T| = \lfloor N_n / 2 \rfloor$, which forces $S = T$ and $f_S = f_T$. Therefore $\mathcal{F}_n$ is an antichain in $\FD(n)$ of cardinality $|\mathcal{F}_n| = \binom{N_n}{\lfloor N_n / 2 \rfloor}$. Combining this with Proposition~\ref{prop:invariance} and Lemma~\ref{lem:binom-est}(1) establishes~\eqref{eq:iter-cb}.

For $n \ge 2$, Lemma~\ref{lem:binom-est}(2) gives $N_n \ge 2^{\lfloor n/2 \rfloor} = 2 \cdot 2^{\lfloor n/2 \rfloor - 1}$, so $\lfloor N_n / 2 \rfloor \ge 2^{\lfloor n/2 \rfloor - 1}$. Applying Lemma~\ref{lem:binom-est}(2) a second time with $M = N_n$ yields
\[
  w(\FD(n)) \;\ge\; \binom{N_n}{\lfloor N_n / 2 \rfloor} \;\ge\; 2^{\lfloor N_n / 2 \rfloor} \;\ge\; 2^{2^{\lfloor n/2 \rfloor - 1}},
\]
proving~\eqref{eq:clean-dexp}.

For the sign classification $\Delta(n) > 0 \iff n \ge 3$: we already proved $\Delta(0) = \Delta(1) = \Delta(2) = 0$ and $\Delta(3) = 1 > 0$. For $n \ge 4$, we have $N_n = \binom{n}{\lfloor n/2 \rfloor} \ge \binom{n}{2} = \frac{n(n - 1)}{2} \ge 6 \ge 4$, so Lemma~\ref{lem:binom-est}(3) applied to $M = N_n$ gives $w(\FD(n)) \ge \binom{N_n}{\lfloor N_n / 2 \rfloor} > N_n$, i.e., $\Delta(n) > 0$.

Finally, let $n_0, B, C, D \in \mathbb{N}$ be arbitrary. Set $K = B + 1$, $M = C + D + 2$, $m = \max(4, n_0, 2K + M + 1)$, and $n = 2(m + 1) \ge n_0$. Since $m \ge 4$, an easy induction gives $2^m \ge m^2$. Moreover, $m \ge 2K + M + 1$ implies
\[
  m^2 \;\ge\; (2K + 1)m \;=\; 2Km + m \;\ge\; 2Km + 2K + M + 1 \;=\; Kn + M + 1.
\]
Since $\lfloor n/2 \rfloor - 1 = m$,~\eqref{eq:clean-dexp} yields
\[
  w(\FD(n)) \;\ge\; 2^{2^m} \;\ge\; 2^{(B + 1)n + C + D + 3} \;=\; 2^n \cdot (2^B)^n \cdot 2^{C + D + 3}.
\]
Set $X = 2^n \ge 1$, $Y = (2^B)^n \ge \max(1, B^n)$ (as $B < 2^B$), and $Z = 2^{C + D + 3} \ge C + D + 3$. Since $X, Y \ge 1$, we have
\[
  X Y Z \;\ge\; X Y (1 + C + (D + 2)) \;=\; X Y + C X Y + (D + 2) X Y \;\ge\; 2^n + C B^n + D + 2.
\]
Because $N_n = \binom{n}{\lfloor n/2 \rfloor} \le 2^n$, subtracting $N_n$ gives $\Delta(n) = \Delta_0(n) = w(\FD(n)) - N_n \ge C B^n + D + 2 > C B^n + D$. Taking $B = 2$ (using $\lfloor \log_2 n \rfloor \le n \le 2^n$) refutes $\Delta(n) = O(\log n)$, and taking $B = 2^d$ (using $n^d \le (2^n)^d = (2^d)^n$) refutes $\Delta(n) = O(n^d)$.
\end{proof}

\begin{remark}[True Asymptotic Order of $\log_2 w(\FD(n))$]
Trivially $w(\FD(n)) \le |\FD(n)| = M(n)$. By Hansel's theorem~\cite{Hansel1966}, $M(n) \le 3^{\binom{n}{\lfloor n/2 \rfloor}}$, and by Kleitman's theorem~\cite{Kleitman1969} (subsequently sharpened by Kleitman--Markowsky~\cite{KleitmanMarkowsky1975} and Korshunov~\cite{Korshunov1981}), $\log_2 M(n) = (1 + o(1))\binom{n}{\lfloor n/2 \rfloor}$. Combining this upper bound with the lower bound~\eqref{eq:iter-cb} yields
\[
  \log_2 w(\FD(n)) = (1 + o(1)) \binom{n}{\lfloor n/2 \rfloor} \qquad (n \to \infty).
\]
Thus Conjecture~00000008532 misidentifies the scale of $w(\FD(n))$ by an exponential function: $\binom{n}{\lfloor n/2 \rfloor}$ governs $\log_2 w(\FD(n))$, whereas $w(\FD(n))$ itself is doubly exponential in $n$ (already reaching $\binom{252}{126} \approx 4.58 \times 10^{74}$ on the middle slice alone at $n = 10$, compared with $\binom{10}{5} = 252$).
\end{remark}

\section{Formal Verification in Lean 4 and Computational Suite}

Every structural lemma, small-$n$ evaluation, and asymptotic refutation in Sections~2--4 has been formalized from first principles in Core Lean~4.19.0 (\texttt{leanprover/lean4:v4.19.0}) inside \texttt{lean/}, compiled with \texttt{-DwarningAsError=true}, zero external dependencies, zero \texttt{sorry} or \texttt{admit}, zero \texttt{native\_decide}, and zero custom axioms. Table~\ref{tab:lean-map} records the correspondence between the mathematical results of this paper and the machine-checked declarations in \texttt{SpernerFD/}.

\begin{table}[ht]
\centering
\caption{Correspondence between mathematical results and Core Lean~4.19.0 declarations.}
\label{tab:lean-map}
\smallskip
{\small
\begin{tabular}{@{}lll@{}}
\toprule
Result & Module & Lean~4 Declaration \\
\midrule
Prop.~\ref{prop:universal} (Universal property) & \texttt{Lattice.lean} & \texttt{evalBounded\_isHom}, \texttt{evalBounded\_unique} \\
Prop.~\ref{prop:invariance} ($w(\FD) = w(\FDz)$) & \texttt{Antichain.lean} & \texttt{spernerFD0\_eq\_spernerFD} \\
Prop.~\ref{prop:conj-lb} ($w(\FD(n)) \ge N_n$) & \texttt{Antichain.lean} & \texttt{centralBinom\_le\_spernerFD} \\
Thm.~\ref{thm:small-and-sign} ($n = 0,1,2,3$) & \texttt{SmallN.lean} & \texttt{spernerFD\_zero} \dots\ \texttt{spernerFD\_three} \\
Thm.~\ref{thm:small-and-sign} ($w(\FD(4)) \ge 24$) & \texttt{SmallN.lean} & \texttt{spernerFD\_four\_ge\_twenty\_four} \\
Lem.~\ref{lem:mid-slice} (Middle-slice antichain) & \texttt{Main.lean} & \texttt{isAntichain\_midSliceAntichain} \\
Lem.~\ref{lem:binom-est}(1)--(3) (Binomial bounds) & \texttt{Basic.lean} & \texttt{centralBinom\_ge\_two\_pow\_half} \\
Thm.~\ref{thm:small-and-sign} ($\Delta(n) > 0 \iff n \ge 3$) & \texttt{Main.lean} & \texttt{spernerCorrectionFD\_pos\_iff} \\
Sign refutation ($n = 1, 2$) & \texttt{Main.lean} & \texttt{not\_always\_positive\_FD\_pos\_n} \\
Thm.~\ref{thm:double-exp} (Iterated binomial LB) & \texttt{Main.lean} & \texttt{spernerFD\_ge\_iterated\_centralBinom} \\
Thm.~\ref{thm:double-exp} (Doubly-exponential LB) & \texttt{Main.lean} & \texttt{spernerFD\_ge\_double\_exp} \\
Thm.~\ref{thm:double-exp} ($\Delta(n) > C B^n + D$) & \texttt{Main.lean} & \texttt{correction\_exceeds\_any\_exponential} \\
Refutation of $O(\log n)$ ($\FD$) & \texttt{Main.lean} & \texttt{not\_logarithmic\_order\_FD} \\
Refutation of $O(\log n)$ ($\FDz$) & \texttt{Main.lean} & \texttt{not\_logarithmic\_order\_FD0} \\
Master refutation ($\FD$) & \texttt{Main.lean} & \texttt{conjecture\_00000008532\_false\_FD} \\
Master refutation ($\FDz$) & \texttt{Main.lean} & \texttt{conjecture\_00000008532\_false\_FD0} \\
\bottomrule
\end{tabular}
}
\end{table}

In addition, the standalone Python~3.12 script \texttt{verify.py} independently generates $\FD(n)$ and $\FDz(n)$ for $0 \le n \le 5$, certifies all exact widths in Table~\ref{tab:small-n} via Hopcroft--Karp matching and K\H{o}nig--Dilworth dual certificates, and verifies the middle-slice antichains and doubly-exponential bounds through $n = 10$.

\clearpage
\begin{thebibliography}{9}

\bibitem{Birkhoff1937}
G.~Birkhoff, \emph{Rings of sets}, Duke Math.\ J.\ \textbf{3} (1937), no.~3, 443--454.

\bibitem{Dedekind1897}
R.~Dedekind, \emph{\"{U}ber Zerlegungen von Zahlen durch ihre gr\"{o}ssten gemeinsamen Theiler}, Festschrift der Technischen Hochschule zu Braunschweig (1897), 1--40.

\bibitem{Dilworth1950}
R.~P.~Dilworth, \emph{A decomposition theorem for partially ordered sets}, Ann.\ of Math.\ (2) \textbf{51} (1950), 161--166.

\bibitem{Engel1997}
K.~Engel, \emph{Sperner Theory}, Encyclopedia of Mathematics and its Applications, vol.~65, Cambridge Univ.\ Press, Cambridge, 1997.

\bibitem{Hansel1966}
G.~Hansel, \emph{Sur le nombre des fonctions bool\'{e}ennes monotones de $n$ variables}, C.\ R.\ Acad.\ Sci.\ Paris S\'{e}r.\ A-B \textbf{262} (1966), A1088--A1090.

\bibitem{Kleitman1969}
D.~Kleitman, \emph{On Dedekind's problem: The number of monotone Boolean functions}, Proc.\ Amer.\ Math.\ Soc.\ \textbf{21} (1969), 677--682.

\bibitem{KleitmanMarkowsky1975}
D.~Kleitman and G.~Markowsky, \emph{On Dedekind's problem: The number of isotone Boolean functions.\ II}, Trans.\ Amer.\ Math.\ Soc.\ \textbf{213} (1975), 373--390.

\bibitem{Korshunov1981}
A.~D.~Korshunov, \emph{The number of monotone Boolean functions}, Problemy Kibernet.\ \textbf{38} (1981), 5--108.

\bibitem{Sperner1928}
E.~Sperner, \emph{Ein Satz \"{u}ber Untermengen einer endlichen Menge}, Math.\ Z.\ \textbf{27} (1928), 544--548.

\end{thebibliography}

\end{document}
